% 本文件由 LaTeX 排版 Agent 根据有效 translation.json 自动生成。
% author=Augustine of Hippo; work_id=1308; title=论节制

\chapter{\WorkTitle{论节制}}

% structure: p0001 source=h1 target=omitted reason=page-title-already-rendered-by-parent

圣\Person{奥斯定}在他的\WorkTitle{书信}第231号\WorkTitle{致达里乌斯伯爵书}中提及他的作品\WorkTitle{论节制}。\EditorialNote{参见本版本第1卷第584页。——P.S.} 波西迪乌斯在\WorkTitle{指示录}第10章中提到了它，伯达或弗洛吕斯的\WorkTitle{集录}以及欧吉皮乌斯也引用过它。因此\Person{伊拉斯谟}根据文体将它归于\Person{雨果}是错误的，它并不像早期讲道录的文体。它显然是一篇讲道，因此很可能未被\WorkTitle{修正录}提及。摩尼教异端受到谴责，方式与其早期著作相似。\EditorialNote{摘自本笃会版本第六卷}\par

1. 要完全恰当而相称地论述灵魂的德行，即所谓的\OriginalTerm{节制}{Continence}，是困难的；但是，这位此德行的伟大恩赐者的天主，将帮助我们微小者承担如此重大的重担。因为祂赐予祂忠信者节制，使他们保持节制，祂自己也赐予祂的仆人们有关节制的话语。最后，我们将论述如此重大的主题，首先我们要说并证明节制是天主的恩赐。我们在\BibleBook{智慧}{篇}上读到：\BibleQuote{除非天主赐与，没有人能节制。}但主论到那更大更光荣的\OriginalTerm{节制}{Continence}本身，即不结婚的节制，说：\BibleQuote{不是所有人都能领受这话，只有那些得了恩赐的人。}而且，既然婚姻的贞洁本身也只有通过避免非法交往的节制才能保持，宗徒在论到两种生活，即婚姻生活和非婚姻生活时，宣称两者都是天主的恩赐，说：\BibleQuote{我愿意众人像我一样，但每人各有自己从天主得来的恩赐；这人这样，那人那样。}\par

2. 而且，唯恐人们以为只应在肉体下体情欲方面从主那里盼望必要的节制，\BibleBook{圣}{咏}中也唱道：\BibleQuote{上主，求你在我的口边派个守卫，在我唇前设立一道节制的门户。}但在神圣言语的这个见证中，如果我们按应当理解的方式理解\Quote{口}，我们就会看到那里设立的节制是天主多么大的恩赐。因为，控制身体的口，防止任何不恰当的东西通过声音爆发出来，这是小事。因为内里还有心灵的口，那说这些话并为我们写下这些话的人，正是向主祈求为他设立节制的守卫和门户。因为有许多事我们不用身体的口说出，却在心中大声呼喊；但身体的口中不会发出任何心中沉默的事物的话语。因此，不从那里流出的东西，不会发出声音；但从那里流出的东西，如果是邪恶的，即使不动舌头，也玷污灵魂。因此，节制必须设立在即使沉默之人的良心也在说话的地方。因为通过节制的门户，可以防止那即使肉体嘴唇紧闭时也污染思想者生命的东西从那里流出。\par

3. 最后，为了更清楚地表明他用这些话所指的内在的口，他在说了\BibleQuote{上主，求你在我的口边派个守卫，在我唇前设立一道节制的门户}之后，紧接着就说：\BibleQuote{不要使我的心偏向邪恶的言语。}心的偏向是什么，岂不就是同意？因为，谁若没有心的偏向，即没有同意所遇到的各事物的刺激，他还没有说话。但是，如果他同意了，他已经在心中说了，即使没有用口发出声音；即使没有用手或身体的任何部分去做，他已经在思想上决定了要做的事：在天主法律面前有罪，虽然对人感官隐藏；言语已在心中说出，行动却未通过身体完成。但是，他绝不会在行动中移动肢体，如果那行动的开端没有事先在言语中发生。因为经上写的不是谎言：\BibleQuote{万事的开端是一句话。}事实上，人们做了许多事，嘴唇闭着，舌头静默，声音受控制；但他们并没有通过身体的行为做出任何他们先前没有在心中说出的事。因此，既然有许多内心的言语所犯的罪并不表现为外在的行为，而外在行为没有不经由内在言语先行的，那么，如果我们围绕内在的嘴唇设立了节制的门户，两者都会有纯洁的无辜。\par

4. 为此，主自己也曾亲口说：\BibleQuote{你们要清洁里面，外面也就洁净了。}此外，在另一处，当祂驳斥犹太人愚昧的言论——他们因门徒不洗手吃饭而毁谤门徒时，祂说：\BibleQuote{不是入口的使人污秽，而是出口的才使人污秽。}若将这整句话理解为肉躯的口，那便是荒谬的。因为食物不使人污秽，呕吐物也不会使人污秽。诚然，食物进入口，呕吐物从口而出。但无疑，前一句论及肉躯的口，祂说：\BibleQuote{不是入口的使人污秽}，后一句却论及内心的口，祂说：\BibleQuote{但出口的，才使人污秽。}最后，当宗徒伯多禄求祂解释这如同比喻的话时，祂回答说：\BibleQuote{你们还不明白吗？岂不知凡入口的，是进到肚腹，又落到茅厕里吗？}这里我们显然看到食物进入的肉躯之口。但祂接下来补充的话，为使我们能认出内心的口，我们愚钝的心本无法跟上，若非真理甚至屈尊与迟钝者同行。因为祂说：\BibleQuote{但从口里出来的，是从心里发出来的。}好像祂在说：当你听见\BibleQuote{从口里}，要理解为\BibleQuote{从心里}。我两者都说，但以一者说明另一者。内在的人有一个内在的口，这口由内在的耳朵分辨；从这个口出来的，是从心里出来的，它们使人污秽。然后，祂放下了\BibleQuote{口}这个也可指身体的词，更明确地揭示了祂所说的：\BibleQuote{因为从心里发出来的，有恶念、凶杀、奸淫、邪淫、偷盗、妄证、毁谤；这些都是使人污秽的。}在这些恶事中，没有一件不能由身体肢体实行，但恶念先已发生，并使人污秽，即使有什么阻碍了罪孽的恶行跟进。因为若因未得权能，手未杀人，难道凶手的心因此就洁净无罪吗？又或者，若一个贞洁的女子不愿与淫荡者通奸，那人因此便未在心中与她行奸吗？又或者，若妓女不在妓院，寻找她的那个人因此便未在心中行邪淫吗？又或者，若想以谎言伤害邻人的人缺乏时机和地点，他便已未用内心的口作妄证吗？又或者，若有人畏惧人，不敢用肉舌大声说出毁谤，心中说\BibleQuote{没有天主}，他便在这罪上无罪吗？同样，人所行的一切其他恶行，凡无身体动作、不为身体感觉所知的，都有其隐秘的罪犯，他们仅凭意念中的同意，即内心的恶言，便已受玷污。圣咏作者害怕自己的心陷入此境，便求主将节欲之门（\OriginalTerm{节欲之门}{door of Continence}）安设于这内心的口唇四周，以约束内心，使它不陷入恶言：约束它，不容许意念进展为同意；如此，按宗徒的训诫，罪便不在我们必死的身体内为王，我们也不将肢体作为不义的武器献给罪。那些因不得权能而不将肢体转向罪恶的人，显然远离了此训诫的实践；若权能临到，他们便立即以肢体动作如同武器那样显明，谁在他们内里为王。因此，就自身而言，他们将肢体作为不义的武器献给罪，因为他们所愿的正是如此，之所以未献，只是因为不能。\par

5. 因此，那因贞洁约束生育部分而最首要、最恰当地被称为节欲（\OriginalTerm{节欲}{Continence}）的，若我们在心中持守了我们刚才所论及的那更高的节欲，便不会因任何过犯而受到侵犯。为此，主说了\BibleQuote{因为从心里发出来的有恶念}之后，接着列举了属于恶念的：\BibleQuote{凶杀、奸淫}等等。祂并未列举全部，而是列举了某些作为例子，教导我们也要理解其他的。其中没有一件能在恶念未先行的情况下发生，恶念先在内在预备了那外在实行的事，并从内心的口发出，已使人污秽，尽管因未得权能，而未通过身体肢体在外实行。因此，当节欲之门安设在内心的口——一切使人污秽之物由此而出——上时，若不容许任何此类事物由此而出，随之而来的便是良心可欢欣的纯洁，尽管尚未达到那不再与私欲争战的完美境地。但现在，只要\BibleQuote{肉欲相反圣神，圣神相反肉欲}，我们能做到不从我们内里所感觉的恶事上同意，便已足够。但当同意发生时，便从内心的口发出那使人污秽的事物。而若凭节欲保持了不同意，肉欲之恶——圣神之欲与之相争——便不能为害。\par

6. 但善战是一回事，这在目前就是抵抗死亡的争战时；没有仇敌则是另一回事，那将在以后，当死亡，\BibleQuote{最后的敌人}，被消灭时。因为\OriginalTerm{节制}{Continence}本身，当它克制\OriginalTerm{情欲}{lusts}时，既寻求我们所向往的不朽之\OriginalTerm{善}{good}，又拒绝我们在此终有一死中所斗争的\OriginalTerm{恶}{evil}。它确实是前者的爱慕者和瞻仰者，却是后者的仇敌和见证者：既追求合宜之事，又逃避不合宜之事。如果我们没有与合宜相悖的意愿，如果邪恶的情欲没有反对我们的善意，那么节制肯定不必费力克制情欲了。\Person{宗徒}高声喊道：\BibleQuote{我知道，}他说，\BibleQuote{那住在我内，就是在我肉身内，并不存在善。因为愿意的善我固能，但实行善我却找不到。}因为如今能行善，在于不赞同邪恶的情欲；但善将完成，当邪恶的情欲本身终止时。同样，这位外邦人的导师也高声喊道：\BibleQuote{我按内在的人喜悦天主的法律，但我看见在我的肢体中另有一条法律，反对我理智的法律。}\par

7. 这种冲突，除了那些为德行而战、为战胜恶习的人以外，没有人亲身体验；除了节制的善以外，没有什么能攻击情欲的恶。但有些人完全不明白天主的法律，不把邪恶的情欲视为仇敌，因悲惨的盲目而做其奴隶，甚至以为满足情欲而非克制情欲就是有福。然而，那些藉着法律而认识情欲的人（\BibleQuote{因为藉着法律乃是对罪的认识}，并且，\BibleQuote{情欲，他说，我若不藉着法律说\InnerQuote{你不可贪恋}，我就不知道}），却在攻击之下被战胜，因为他们活在法律之下，法律命令善但不赐予；他们不活在恩宠之下，恩宠藉着圣神赐下法律所命令的。因此法律进入，为要使过犯增多。禁令增加了情欲，使它不可战胜；这样也就有了过犯，没有法律时虽有过犯，却没有过犯。\BibleQuote{因为哪里没有法律，哪里就没有过犯。}这样，法律因不帮助恩宠而禁止罪恶，反而成了罪恶的力量；故\Person{宗徒}说：\BibleQuote{法律是罪恶的力量。}这并不奇怪，因为人的软弱甚至从良善的法律中给邪恶增添了力量，同时却相信能凭自己的能力遵守法律本身。他们认识不到天主在软弱者中赐下的义德，想要建立自己的义德（软弱者所缺乏的），因此没有服从天主的义德，成了可弃绝而骄傲的人。但如果法律作为\OriginalTerm{启蒙师}{schoolmaster}，领受过犯者走向恩宠，似乎为此目的而更严重地伤害他，使他渴求\OriginalTerm{良医}{Physician}；那么，主就赐予一种产生善的甘甜，以对抗情欲所胜的毒害之甜，使节制的甘甜更令人喜悦，并且\BibleQuote{我们的地就要出产她的果实}，使那藉着主的帮助而战胜罪恶的士兵得以饱饫。\par

8. 宗徒的号角以这样的声音激励这些士兵投入战斗：\BibleQuote{所以，不要让罪恶在你们终有一死的身体上为王，顺从它的情欲；也不要把你们的肢体献给罪恶作不义的武器；但要献给你们自己给天主，如同从死者中复活的人，并把你们的肢体献给天主作正义的武器。因为罪恶将不得统治你们，因为你们不在法律之下，而在恩宠之下。}在另一处：\BibleQuote{所以，弟兄们，我们并不是欠肉体的债，去顺从肉体生活。如果你们顺从肉体生活，必要死亡；但如果你们依赖圣神，治死肉体的行为，必要生活。因为凡被天主的圣神所引导的，都是天主的儿子。}这就是当前的任务，只要我们在恩宠之下的此世生命持续不断，就是要让罪恶，即情欲的罪恶（他在这里以罪恶之名称呼它），不在我们终有一死的身体内为王。但如果服从它的欲望，它就显出为王。因此，我们里面有情欲的罪恶，它不得为王；有它的欲望，我们不得服从，以免因服从它而在我们身上为王。因此，不要让情欲夺取我们的肢体，而让节制为自己要求它们；使它们成为献给天主的正义的武器，而不是献给罪恶的不义的武器；这样，罪恶就不能统治我们。因为我们不在法律之下，法律确实命令善却不赐予；我们在恩宠之下，恩宠使我们爱法律所命令的，因而能统治自由者。\par

9. 并且，当他劝诫我们，不要随从肉身生活，免得死亡，但要依靠圣神致死肉体的作为，好使我们生活时；那吹响的号角确实显示了我们所参与的战争，激发我们奋力交战，将敌人置于死地，以免我们被他们所杀。然而，它已足够清楚地指明了那些敌人是谁。因为那些就是它愿意我们置之于死地的，即肉体的作为。因为经上说：\BibleQuote{但如果你们依靠圣神致死肉体的作为，你们必要生活。}为了让我们知道这些是什么，让我们听一听同样写给迦拉达人的书卷中所说的话：\BibleQuote{但肉体的作为是显露的，即淫乱、不洁、奢侈、崇拜偶像、施行邪法、仇恨、竞争、嫉妒、忿怒、争吵、异端、妒忌、醉酒、宴乐，以及类似的事。关于这些，我预先告诉你们，正如我先前所说过的：行这样事的人，必不能承继天主的国。}因为那场战争，他也同样说明，他应该谈论这些，并用同一把属天和属神的号角呼召基督的士兵去彻底消灭这些敌人。因为他上面曾说过：\BibleQuote{但我告诉你们：在圣神内行走，而不去满足肉体的私欲。因为肉体与圣神相争，圣神与肉体相争；两者彼此敌对，致使你们不能做你们所愿意的。但如果你们被圣神引导，就不在法律之下。}因此，他们既处于恩宠之下，他就愿意他们与肉体的作为进行那场战斗。为了指出这些肉体的作为，他加上了我刚才所提到的话。\BibleQuote{但肉体的作为是显露的，即淫乱、}以及其余，无论是他所提及的，还是他所告诫的，都要当作主要被他所加的\Quote{以及类似的事}来理解。最后，在这场战斗中，针对可以说是肉体的军队，他仿佛调遣了另一支属神的队伍，说：\BibleQuote{但圣神的果实是：爱德、喜乐、平安、忍耐、仁慈、良善、信德、温和、节制；反对这一类的事，无法可制。}他没有说\Quote{反对这些，}以免人们以为只有这些；即使他这样说，我们也当理解一切，凡我们所能想到的同类的善行；但他却说\Quote{反对这一类，}就是说，既包括这些，也包括一切类似的事。然而，在他所提及的善行中，他将节制置于最后（关于节制我们现在已着手论述，并已谈了许多），他愿意它以一种特殊的方式铭刻在我们心中。确实，节制在圣神与肉体相争的情况下大有裨益，因为它在某种程度上将肉体的私欲钉在十字架上。因此，宗徒说了这些话后，立刻加上：\BibleQuote{凡属于耶稣基督的人，已把自己的肉身连同私欲偏情钉在十字架上了。}这就是节制的作为：这样，肉体的作为就被置于死地。但那些肉体的作为却将那些因背离节制而被私欲牵引，去同意行这样事的人置于死地。\par

10. 但为了使我们不背离节制，我们应当特别警惕魔鬼暗示的圈套，以免我们信赖自己的力量。因为\BibleQuote{凡信赖世人的人，是可咒骂的。}除了人之外，还有谁呢？因此，我们不能真正地说，那信赖自己的人不信赖世人。因为\Quote{随从人生活}若不是\Quote{随从肉身生活}又是什么呢？所以，凡受这类暗示诱惑的人，要听一听，并且如果他有任何基督徒的情感，就当战栗。我说，要听一听：\BibleQuote{如果你们随从肉身生活，必要死亡。}\par

11. 但有人会对我说：按人生活是一回事，按肉身生活是另一回事；因为人确实是有理性的受造物，且在他内有理性的灵魂，借此他不同于禽兽；但肉身是人的最低下、最属土的部分，因此按肉身生活是有缺陷的。为此，那按人生活的，当然不是按肉身生活，而是按人之所以为人的那一部分，即按精神的理智而生活，借此他超越禽兽。但这种讨论在哲学学派中或许有些力量；然而我们，为要理解基督的宗徒，应当留意基督宗教的书籍通常是如何表述的；至少我们所有这些以活着就是基督为信念的人，都相信：人被天主圣言所取，绝不可能没有理性的灵魂，正如某些异端者所主张的那样；然而我们读到：\BibleQuote{圣言成了肉身。}这里\Quote{肉身}应理解为别的，不就是人吗？\BibleQuote{凡有血肉的，都要看见天主的救恩。}这能理解为别的，不就是所有的人吗？\BibleQuote{愿万民都到你面前来。}这能理解为别的，不就是所有的人吗？\BibleQuote{你赐给他权柄掌管凡有血肉的。}这能理解为别的，不就是所有的人吗？\BibleQuote{凡有血肉的，都不能因行法律而成义。}这能理解为别的，不就是没有人能成义吗？这一点，同一位宗徒在另一处更明确地承认说：\BibleQuote{人不能因行法律而成义。}他也责备格林多人说：\BibleQuote{你们不还是属血肉的，并且按人而行吗？}在他称他们为属血肉的之后，他没有说你们按肉身而行，而是说按人而行，因为借此他还想指明别的，不就是按肉身吗？因为假如说按肉身而行，即生活，应受责备，而按人而行应受赞美，他就不会以责备的口吻说\BibleQuote{你们按人而行}了。让人认识这责备吧；让他改变意向，逃避毁灭吧。人哪，你听着：不要按人而行，而要按那造人的那位而行。不要从造你的那位那里堕落，甚至到了你自己身上。因为有一位并没有按人而生活的人说：\BibleQuote{我们并不是凭自己足以想什么事，好像是出于自己；我们的充足是出于天主。}请你想想，他是否按人生活，而他又真实地说了这些话？因此，宗徒劝诫人不要按人生活，是为了将人归还于天主。但谁若不按人生活，而按天主生活，他当然也不按自己生活，因为他自己也是一个人。然而当他如此生活时，也被称为按肉身生活；因为当仅提肉身时，人就被理解了，正如我们已经表明的那样；就如同当仅提灵魂时，人也被理解一样：因此经上说：\BibleQuote{人人都应服从上级有权柄的人，}即每一个人；以及\BibleQuote{与雅各伯一同下到埃及的有七十五个灵魂，}即七十五个人。所以，人哪，不要按你自己生活；你曾由此丧亡，但你被寻找了。那么，我再说，不要按你自己生活，人哪；你曾由此丧亡，但你被找到了。当你听到\BibleQuote{如果你们按肉身生活，你们必要死}时，不要指控肉身的本性。因为这话可以这样说，并且最真实地可以说：如果你们按你们自己生活，你们必要死。因为魔鬼没有肉身，然而由于它愿按自己生活，\BibleQuote{它不持守真理。}因此，如果它按自己生活，\BibleQuote{它说谎言时，是出于自己的本性，}这是真理关于它所说的真话。\par

12. 因此，当你听到\BibleQuote{罪不得统治你们}时，不要靠你自己而自信罪不统治你，而要倚靠那一位，一位圣人曾向祂祈祷说：\BibleQuote{请按你的圣言引导我的路径，不要让任何不义辖制我。}因为恐怕在我们听了\BibleQuote{罪不得统治你们}之后，我们却自高自大，将此事归于自己的力量，宗徒立刻看出了这一点，并补充说：\BibleQuote{因为你们不在法律之下，而是在恩宠之下。}因此，是恩宠使罪不统治你。所以，不要靠你自己而自信，免得它因此更加统治你。当我们听到\BibleQuote{如果你们凭圣神治死肉体的作为，你们必要生活}时，我们不要将这如此巨大的美善归于我们自己的灵，仿佛它自己能做这事。为了使我们不抱有那属肉体的想法——（我们的）灵是死的，而非那能治死其它东西的——他立刻补充说：\BibleQuote{因为凡被天主的圣神引导的，都是天主的子女。}因此，为使我们凭我们的灵治死肉体的作为，我们是被天主的圣神引导，祂赐予\OriginalTerm{节德}{Continence}，用以约束、驯服、战胜情欲。\par

在这如此重大的争战中，人在恩宠之下生活，当他获得帮助，善加战斗，便战战兢兢地因主而喜乐。然而，即使是英勇的战士，及那克苦肉身且不可战胜的人，也不免有某些罪的创伤，为了医治这些创伤，他们可以每天说：\Quote{请宽恕我们的罪债。}他们藉着祈祷，以更大的警醒和敏锐，对抗同一邪恶，以及对抗恶魔——诸恶的元首和君王，使他的致命提议不得逞；他更鼓动罪人推卸而非承认自己的罪过，以致那些创伤不但不得医治，反而，即使原本并非致命，也可能被击中要害，造成严重而致命的伤害。因此，在这里需要更为谨慎的节制，以约束人类的骄傲欲望；这欲望使人自满，不愿被人责备，犯罪时不屑被定罪，不肯怀着有益的谦卑承认自己的罪，反而带着致命的傲慢寻找借口。为了约束这种骄傲，那人（他的话我已在上文引述，并尽量加以称赞）向主祈求节制。因为他说完\Quote{上主，求你在我的口前设置守卫，在我唇前围上节制的门。不要使我的心偏向邪言恶语}之后，更清楚地解释为何说这话，便说：\Quote{用来在罪中寻找借口。}还有什么比这些话更邪恶呢？恶人藉此否认自己是恶的，尽管他被指控做了恶事，而这事他无法否认。既然他不能隐藏行为，也不能说它做得好，而且明知这事显然是他所为，他便试图将自己的行为推给别人，仿佛他能移除自己应得的报应。他不愿自己承担责任，反而增加自己的罪责；他藉着推卸而非承认自己的罪过，却不知道他是在推开怜悯而非惩罚。因为在人间的审判者面前，由于他们可能受骗，暂时似乎有些益处，可以藉某种欺骗来清洗所行的错事；但在不能被欺骗的天主面前，我们必须使用真实的罪过宣认，而非欺骗性的辩护。\par

有些人惯于推卸自己的罪过，抱怨自己被命运驱策而犯罪，仿佛星宿命定了此事，是上天先命定这一切而犯罪，好使人后来因行同样的事而犯罪，于是更将自己的罪归咎于运气；他们认为万物都是由偶然意外所推动，却又坚持他们的这种智慧和断言不是出于偶然的轻率，而是出于确定的理性。这是何等的疯狂：既把他们的讨论归于理性，却又使自己的行为受偶然支配！另一些人将自己所做的一切恶事完全归于魔鬼，甚至不愿与他分享罪责，因为一方面他们可能怀疑他是否通过隐秘的暗示诱使他们作恶，另一方面却无法怀疑他们已同意了这些暗示，不论其来源如何。还有些人将自己的辩护扩展到指控天主，他们因天主的审判而痛苦，却因自己的疯狂而亵渎。因为他们引进一种与天主对立的邪恶造反的本体，说天主无法抵抗，除非他将自己本体和本性的一部分掺入那叛逆的实质中，使之被污染和败坏；他们说他们犯罪是在邪恶的本性战胜天主的本性之时。这就是摩尼派最污秽的疯狂，其魔鬼般的计谋被毋庸置疑的真理轻易推翻；这真理承认天主的本性不能被污染和败坏。但人若相信这些人的说法，相信那位至善、无可比拟的天主能够被污染和败坏，那么他们岂不也该被认为有邪恶的污染和败坏吗？\par

15. 还有一些人，为自己的罪过辩解，竟然如此控告天主，说罪过令祂喜悦。他们说：\Quote{如果罪过不令祂喜悦，那么以祂的全能，祂绝对不会容许它们发生。}仿佛天主真的让罪过不受惩罚，甚至对那些祂藉着赦免罪过而免于永罚的人也是如此！诚然，没有人能因应受更重惩罚而得到赦免，除非他已经承受了某种惩罚，无论多轻，都远少于应得的；并且慈悲的满盈如此传达，以至于纪律的公义也不被放弃。因为罪过，即使是看似未受报复的，也有其自身的惩罚相伴随，以致没有人不因他所做的事而受苦：或因痛苦而受折磨，或因盲目而不受苦。因此，正如你所说：\Quote{如果这些事令祂不悦，为什么祂容许它们发生？}我也要说：\Quote{如果它们令祂喜悦，为什么祂又惩罚它们？}这样，正如我承认，除非全能者容许，这些事根本不会发生；同样我也承认，公义者所惩罚的事，我们不应去做。这样，因着不做祂所惩罚的事，我们才配得知，为什么祂容许祂所惩罚的事存在。因为正如经上记载：\Quote{硬食是为成全的人}，那些已有长进的人已经明白，这更属于天主全能的能力，即容许来自自由意志选择的恶存在。的确，祂的全能之善如此伟大，甚至能从恶中生出善来：或藉着赦免，或藉着医治，或藉着使恶适合并转向虔诚者的益处，甚至藉着最公正的报复。因为所有这些都是善的，且最适合良善全能的天主；然而它们只能从恶中产生。所以，还有什么比祂更良善、更全能呢？祂从不制造恶，却甚至能从恶中造出善来？行恶的人向祂呼求：\Quote{宽免我们的罪债}，祂垂听，祂赦免。他们自己的恶伤害了罪人，祂帮助并医治他们的疾病。祂子民的敌人发怒，祂从他们的愤怒中造出殉道者。最后，祂也定那些祂判断为应受惩罚之人的罪；虽然他们承受他们自己的恶，但祂所做的却是善。因为公义的不能不是善的，而且正如罪是不义的，罪过的惩罚却是公义的。\par

16. 但天主并非没有能力造人，使人不能犯罪；祂宁愿造人，使人有能力犯罪，如果他愿意；有能力不犯罪，如果他不愿意；禁止其一，命令另一；这样，对他而言，先是不犯罪的美好功绩，后是公正的奖赏：不再有能力犯罪。因为末后，祂也将使祂的圣人们如此，即完全没有犯罪的能力。诚然，甚至现在祂的众天使也是如此，我们在祂内如此爱他们，以致不为他们中的任何一个担忧，害怕因犯罪而变成魔鬼。而对于任何在这必死生命中的义人，我们不敢这样推测。但我们相信，在永生中，众人都会如此。因为全能的天主，甚至从我们的恶中行善，当祂使我们摆脱一切恶时，祂将赐下何等的美善！关于善用恶，可以更详尽、更微妙地讲论许多；但这并非我们当前论述所承担的任务，我们必须避免过长的篇幅。\par

17. 因此，现在让我们回到我们所说这些的目的上来。我们需要\OriginalTerm{节欲}{Continence}，我们知道这是天主的恩赐，好使我们的心不至于偏向恶言，为罪过作辩解。但有什么罪过，是我们不需要节欲来阻止其犯下的呢？因为正是这节欲，在罪过已经犯下时，阻止它因邪恶的骄傲而被辩护。因此，普遍来说，我们需要节欲，以远离邪恶。但行善似乎属于另一种德行，即义德。神圣的圣咏提醒我们，那里写着：\Quote{远离恶，行善。}但为了什么目的而行善，它随即补充说：\Quote{寻求和平，并追随它。}因为当我们的本性不可分离地依附于它的创造者，我们自身没有任何反对我们自己的东西时，我们将拥有完全的和平。我们的救主自己也愿意我们明白这一点，依我看来，当祂说：\Quote{你们要束紧腰，点上灯}时。束紧腰是什么意思？就是抑制私欲，这是节欲的工作。而点上灯，则是以善工发光发热，这是义德的工作。祂也没有在这里沉默，说明我们做这些事的目的，补充说：\Quote{你们要像等候主人从婚筵回来的人一样。}但当祂到来时，祂将赏报我们，那些守住了自己、远离了私欲所贪恋之事，并行了爱德所命令之事的人，使我们能在祂完全而永恒的平安中为王，没有任何邪恶的争战，而有至高的善的喜乐。\par

所有我们因此，相信永生真天主的人，祂的本性至善且不变，既不作恶，也不受苦，一切美善都从祂而来，甚至那能减少的善，而祂在自己的善——即祂自身——中完全不减少，当我们听见宗徒说：\BibleQuote{在圣神内行走，且不要实行肉体的私欲。因为肉体私欲反对圣神，圣神也反对肉体私欲：它们彼此对立，致使你们不能行你们所愿的。}断不可相信摩尼教人的狂妄，他们认为这里显明两种相反、彼此争斗的本性或本原，一为善的本原，一为恶的本原。实际上这两者都是善的：圣神是善，肉体也是善；由两者组成的人，一个管辖，一个服从，无疑是善，但却是可变化的善，这善若非由不变的善所造，绝不能存在，凡是善，无论大小，都由祂创造；但无论多么小，却由伟大者所造；无论多么大，却绝不能与造物主的伟大相比。但在人的这个本性中，它是善的，由善者所造并妥善安排，如今却有战争，因为尚未痊愈。让疾病得医治，便有和平。但那是罪过应当承受的疾病，而非本性固有的。这罪过确实通过重生的洗涤，天主的恩宠已宽免了信徒；但在同一位医师的手中，本性仍与疾病争斗。但在这样的争斗中，胜利将是完全的痊愈；而且那痊愈不是暂时的，而是永远的：其中不仅这疾病要终结，而且之后也不再有疾病。因此，义人对自己的灵魂说：\BibleQuote{我的灵魂，请赞颂上主，不要忘记祂的一切恩赐：祂赦免你的一切罪愆，医治你的一切疾病。}祂赦免我们的罪愆，当祂赦免罪过；祂医治疾病，当祂抑制邪恶的私欲。祂通过赐予赦免而成为罪愆的仁慈者；祂通过赐予节制而医治疾病。第一件事在圣洗中为告解者成就；第二件事在争斗中为奋战者成就；在其中藉祂的帮助，我们要战胜我们的疾病。甚至现在第一件事也成就，当我们被垂听，说：\BibleQuote{赦免我们的罪债}；而第二件事，当我们被垂听，说：\BibleQuote{不要让我们陷入诱惑。}雅各伯宗徒说，\BibleQuote{因为每个人受诱惑，被自己的私欲所牵引、所诱惑。}为对抗这罪过，我们从祂那里寻求药物的帮助，祂能医治所有这类的疾病，不是通过挪去与我们本性相异之物，而是通过更新我们自己的本性。因此上述宗徒不说\BibleQuote{每个人受诱惑}由于私欲，而是加上\BibleQuote{自己的}：使听见的人明白，他应如何呼号：\BibleQuote{我说：主，求祢怜悯我，医治我的灵魂，因为我得罪了祢。}因为如果它没有因犯罪而败坏自己，使得自己的肉体反对它，即它自身反对自己，在那方面它在肉体内患病，它就不需要医治了。\par

因为肉体不通过灵魂便无所贪求，但说肉体反对心神，当灵魂以肉体的私欲与心神争斗时。这整个是我们：而肉体本身，当灵魂离开时就死去，是我们最低的部分，不是当作我们应逃避之物而丢弃，而是当作我们应再领受之物而放下，领受之后，永不再离弃。但\BibleQuote{播下的是属生灵的身体，复活起来的是属神的身体。}那时，肉体将不再反对心神贪求任何事，因为连它自己也要被称为属神的，不仅毫无反对，而且毫无肉身营养的需要，永被服从于心神，由基督所复活。因此，现在在我们内彼此对立的这两者，既然我们存在于两者之中，让我们祈祷并努力使它们和好。因为我们不应认为其中一个是仇敌，而是那使肉体反对心神的罪过：这罪过得医治后，其本身将止息，两者皆安然无恙，两者间亦无争斗。让我们听宗徒说：\BibleQuote{我知道，在我内，即在我的肉体内，没有善住着。}他确实这样说；在一件善事中的肉体的罪过，不是善；这罪过止息后，它仍是肉体，但不再是腐败或有罪的肉体。然而这属于我们的本性，同一位导师也显示，他先说\BibleQuote{我知道，在我内}，为了解释此句，他加上\BibleQuote{即在我的肉体内，没有善住着。}因此他说他的肉体是他自己。那么它不是我们的仇敌：当抵抗它的罪过时，它本身被爱，因为它被照顾；\BibleQuote{因为从来没有人憎恨自己的肉身}，正如宗徒自己所说。在另一处他说：\BibleQuote{这样看来，我本人是以理智服事天主的法律，而以肉体服事罪的法律。}有耳可听的，让他们听吧。\BibleQuote{这样看来，我本人}：我以理智，我以肉体，但\BibleQuote{以理智服事天主的法律，以肉体服事罪的法律。}如何\BibleQuote{以肉体服事罪的法律}？难道是同意肉体的私欲？绝不可以！而是在其中有不想要的欲望的冲动，却仍有。因不同意它们，他以理智服事天主的法律，并保守其肢体不作罪的武器。\par

20. 因此，在我们内有着邪恶的欲望；若不予以同意，我们便不活得邪恶：在我们内有着罪的私欲；若不予以服从，我们便不完全邪恶，但因为拥有它们，我们也尚未完全为善。宗徒同时指出了这两点：既没有在此处成全善，因为邪恶如此被贪求；也没有在此处成全恶，因为这类私欲未被服从。他确实指出了一方面，当他说：\BibleQuote{愿意善是与我同在，但成全善却不是}；另一方面，当他说：\BibleQuote{你们应在圣神内行走，而不要成全肉身的私欲。} 因为在前一处，他并没有说行善不与他同在，而是说\Quote{成全}；在此处他也没有说\Quote{没有肉身的私欲}，而是说\Quote{不要成全}。因此，在我们内发生邪恶的私欲，当那不合法的事取悦我们时；但它们不被成全，当邪恶的私欲被那服侍天主法律的理智所抑制时。善也发生，当那错误地取悦的事，因良善的喜悦占优势而不能得逞时。但是善的成全尚未完成，只要肉身服侍罪的律，邪恶的私欲诱惑，并且虽然被抑制，却仍然在活动。因为如果它不动，就没有必要抑制它。在某时也将有善的成全，那时恶被毁灭：前者将是最高的，后者将不复存在。如果我们认为在此有死的状态下可以盼望这个，我们就受了欺骗。因为那将在死亡不再存在的时候；将在那里，那里有永生。因为在那个世界，那个国度，将有至善，没有邪恶：那时，将会有对智慧的最高爱慕，而没有节制的劳苦。因此，肉身并非邪恶，如果它没有邪恶，即没有过犯；由于这过犯，人变得有过错，不是被人弄坏，而是自己弄坏的。因为无论哪一部分，即灵魂和身体，都由善的天主造成善的，但人自己制造了那使他变成恶的恶。借着宽恕，人已经从这恶的罪咎中得到释放，为了不使他以为自己所做的是轻事，他仍然借着节制与自己的过犯争战。但在那将来的平安中统治的人中，但愿没有任何过犯；因为在这战争的状态中，在进步的人身上，每天减少的不仅是罪，而且正是私欲本身；我们不与其同意而争战，若同意就犯罪。\par

21. 因此，肉身相反圣神，在我们肉身内没有善住着，我们肢体中的法律反对理智的法律，这并不是由相反原则所引起的两种本性的混合，而是由于罪的报应所造成的一个分裂反对自己。在亚当内，在我们之前，当那本性听从并追随了它的欺骗者，而轻视并冒犯了它的造物主时，我们并不是这样：这并非被造的人的先前生命，而是被判罪的人的后来的惩罚。当人从这定罪中借着恩宠、通过耶稣基督得到释放时，被释放的人与他们的惩罚争战，尚未得到完全的救恩，但已得到救恩的抵押；但若未得释放，他们便因罪而有罪，且被卷入惩罚之中。但在今世之后，对于有罪者，他们的罪行将永远留下惩罚；对于被释放者，将不再永远留下罪行或惩罚：但善的实体，即灵魂和肉身，将继续永远存在，这些是天主——祂是善的、不能改变的——所创造的，祂造它们为善，虽然是可变的。但它们将继续存在，已被转变为更好的，从此再不转变为更坏的：一切邪恶都被彻底毁灭，无论是人不义地所做的，还是他正义地所承受的。这两种邪恶完全消亡，其一是先行的不义，其二是后来的不幸；人的意志将正直而无任何堕落。在那里，将清楚明了地呈现在所有人面前，现在许多信徒相信、少数人理解的事：恶不是实体；而是如同身体上的创伤，在一个使自己变得有过错的实体中，当疾病开始时，它开始存在，当痊愈完成时，它就不复存在。因此，一切邪恶已从我们兴起，并在我们内被毁灭，我们的善也被增长并完善到最幸福的不朽和不死的崇高，那么我们的两个实体将会是怎样的呢？因为如今，在这腐朽与必死中，正如\BibleQuote{那腐朽的身体压住灵魂；} 并且宗徒所说的：\BibleQuote{身体因罪而死亡；} 然而同一位宗徒这样为我们的肉身作证（即我们最低的、属地的部分）说道，如我刚才所提到的：\BibleQuote{从来没有人憎恨过自己的肉身。} 他紧接着又说：\BibleQuote{反而保养爱惜它，如同基督对待教会一样。}\par

22. 因此，我所说的不是有何等错误，而是何等的全然疯狂，竟然使摩尼派把我们肉身的由来归诸于某个我不知其名的、传说中所谓的\Quote{黑暗种族}，他们竟说这种族自有其本性，无始而常恶；而真正的导师却劝人以其自身肉身的模式去爱自己的妻子，并且也以基督和教会的模式来劝他们做同一件事。最后，我们必须忆起宗徒书信中与此事密切相关的那整段经文。\BibleQuote{丈夫们，}他说，\BibleQuote{你们要爱自己的妻子，如同基督也爱了教会，并为她舍了自己，为圣化她，藉着话中之水的洗涤而洁净她，为把她呈献给自己，成为光荣的教会，没有斑点、皱纹或任何这类东西，但使她圣洁无瑕。同样，}他说，\BibleQuote{丈夫也应当爱自己的妻子，如同爱自己的身体。爱自己妻子的，就是爱自己。}他接着说，我们已经提及的那句：\BibleQuote{从来没有人恨恶自己的肉身，而是养之育之，如同基督对待教会一样。}对这一切，最不洁不敬的疯狂又作何回应？你们，摩尼派，对这些话作何回应？你们竟想好像从宗徒书信中给我们带来两种无始的本性，一善一恶，却不愿听从宗徒书信，好从那种亵渎的乖戾中纠正你们？你们读到\BibleQuote{肉身逆于心神}和\BibleQuote{我肉身内无善住}；同样也请读到\BibleQuote{从来没有人恨恶自己的肉身，而是养之育之，如同基督对待教会一样。}你们读到\BibleQuote{我发觉在我肢体内另有一条法律，与我理智的法律相敌对}；同样也请读到\BibleQuote{正如基督爱了教会，男人也当爱自己的妻子，如同爱自己的身体。}你们不要在前面的圣经见证上狡猾，而在后面的见证上耳聋，这样你们在两处都会正确。因为，如果你们按正理接受后者，你们就会努力按真理理解前者。\par

23. 宗徒让我们知道三种结合：基督与教会、丈夫与妻子、心神与肉身。在这些结合中，前者为后者的益处着想，后者服事前者。当一切事物中，本应处于优等地位的被设立，本应适当服从的也服从时，它们便都遵守秩序之美。丈夫与妻子领受了命令与模式，他们应如何彼此相待。命令是：\BibleQuote{妻子要顺服自己的丈夫，如同顺服主，因为丈夫是妻子的头}；以及\BibleQuote{丈夫们，要爱你们的妻子。}但模式也给了妻子从教会，给了丈夫从基督：\BibleQuote{正如教会顺服基督，同样妻子也凡事顺服自己的丈夫。}同样，在命令丈夫爱自己的妻子之后，他又加上一个模式：\BibleQuote{正如基督爱了教会。}但丈夫也从一件较低的事上受到劝勉，即从自己的身体：不仅是从更高的，即从他们的主。因为他不仅说\BibleQuote{丈夫们，要爱你们的妻子，如同基督也爱了教会}，这是从更高的；他也说\BibleQuote{丈夫应当爱自己的妻子，如同爱自己的身体}，这是从较低的：因为高和低都是好的。然而，女人并没有从身体或肉身上得到模式，要像肉身顺服心神那样顺服丈夫；或者宗徒希望我们按推论理解他所省略的陈述；或者或许因为在这个有死且有病的生命状态中，肉身逆着心神，所以他不想从肉身给女人设立顺服的模式。但他却给男人这个理由，因为尽管心神逆着肉身，但即使在此情况下它也为肉身的益处着想：不像肉身那样逆着心神，这种对立既不为心神的益处着想，也不为自身着想。然而，善良的心神不会为肉身的益处着想，无论是通过预先思想养育和爱护其本性，还是通过节制抵抗其过犯，若非每种实体藉着其秩序的适宜性，都显明天主是两者的创造者。因此，你们既以真正的疯狂自称为基督徒，又如此乖戾地反对基督徒圣经，闭着眼睛，不如说是瞎了眼睛，断言基督曾以虚假的肉身显现于世人，而教会中灵魂属于基督，身体属于魔鬼，男性与女性的性别是魔鬼的作为，不是天主的作为，肉身与心神结合如同恶实体与善实体结合，这究竟是为何？\par

24. 如果我们从\OriginalTerm{宗徒书信}{Apostolic Epistles}中所提及的，在你看来尚不足以回答，那么，如果你有耳朵，请再听听别的。极其疯狂的\OriginalTerm{摩尼教徒}{Manichæan}对\OriginalTerm{基督的肉身}{Flesh of Christ}说了什么？他说那肉身不是真实的，而是虚假的。对此，蒙福的宗徒又说了什么？\BibleQuote{你务要记着，达味的后裔耶稣基督从死者中复活了，这符合我所传的福音。} 而且基督耶稣亲自说：\BibleQuote{你们摸摸看，并看看：鬼神是没有骨肉的，如同你们看我一样。} 那么，在他们的教义中，声称基督的肉身是虚假的，哪里有真理呢？在如此巨大的谎言中，基督身上怎会没有邪恶呢？因为，对过于洁净的人来说，真实的肉身是邪恶，而虚假的肉身代替真实的，就不是邪恶吗？出自达味后裔的真实肉身是邪恶，而说\BibleQuote{你们摸摸看，并看看：鬼神是没有骨肉的，如同你们看我一样}的虚假舌头，就不是邪恶吗？对于教会，这个以致命的错误欺骗世人之人说了什么？他说，从灵魂方面，教会属于基督，从身体方面，则属于魔鬼？对此，在信德和真理上的\OriginalTerm{外邦人的导师}{Teacher of the Gentiles}说了什么？他说：\BibleQuote{你们不知道你们的身体是基督的肢体吗？} 对于男女的性别，这个丧亡之子说了什么？他说，两性都不是出于天主，而是出于魔鬼。对此，\OriginalTerm{蒙拣选的器皿}{Vessel of Election}说了什么？他说：\BibleQuote{原来女人是由男人而出，男人也是由女人而生；但一切都出于天主。} 对于肉身，不洁之神通过摩尼教徒说了什么？他说，肉身是邪恶的本体，不是天主的创造，而是仇敌的创造。对此，圣神通过保禄说了什么？他说：\BibleQuote{就如身体只是一个，却有许多肢体，但身体的所有肢体虽说很多，仍是一个身体：基督也是这样。} 稍后又说：\BibleQuote{天主在身体上安排了每个肢体，都如他所愿意的。} 又稍后说：\BibleQuote{天主将更大的光荣赐予那有缺欠的肢体，使身体没有分裂，反倒使各个肢体彼此互相关照：若是一个肢体受苦，所有的肢体都一同受苦；若是一个肢体受尊荣，所有的肢体都一同欢乐。} 如果肉身是邪恶的，为什么连灵魂自己都被劝勉要模仿其肢体的和平？如果肉身是仇敌的创造，为什么连管理身体的灵魂自己都从身体的肢体取得榜样，在彼此之间没有敌意的分裂，好使天主本性上赐予身体的，他们自己也能因恩宠而乐意拥有？有鉴于此，他在致罗马人的书信中写道：\BibleQuote{所以弟兄们！我以天主的仁慈请求你们，献上你们的身体当作生活、圣洁和悦乐天主的祭品。} 如果我们献上出自\Quote{\OriginalTerm{黑暗的民族}{nation of darkness}}的身体，作为生活、圣洁和悦乐天主的祭品，那么我们说黑暗不是光、光不是黑暗，就是没有道理的。\par

25. 但是，他们说，肉身如何能与教会相比拟呢？什么！难道教会相反基督吗？而同一宗徒却说：\BibleQuote{教会服从于基督。}显然，教会服从于基督；因此圣神相反肉身，为使教会在各方面都服从于基督；但肉身相反圣神，是因为教会尚未获得那所应许的圆满平安。为此，教会为救恩的保证而服从于基督，而肉身因疾病的软弱而相反圣神。因为那些他如此说话的人，也正是教会的肢体：\BibleQuote{你们应随从圣神的引导，绝不可满足本性的私欲，因为本性和圣神彼此对立，致使你们不能做你们所愿意的事。}这些话确切地是对教会说的，如果教会不服从基督，圣神就不会藉着节德在它内相反肉身。因此，他们确实能够不满足肉身的私欲，但因着肉身相反圣神，他们不能做他们愿意的事，即连肉身的那些私欲本身也不能有。最后，我们为何不承认在属神的人身上，教会服从基督，而在属血气的人身上，教会却还相反基督呢？难道那些人对基督没有私欲吗？对他们曾说过：\BibleQuote{基督被分裂了吗？}以及：\BibleQuote{我不能对你们讲话如同对属神的人，只能如同对属血气的人。我把你们当作在基督内的婴孩，给你们奶喝，没有给饭食，因为你们当时还吃不下；就是如今，你们还是吃不下；因为你们仍是属血气的。在你们中间有嫉妒和纷争，你们岂不是属血气的吗？}嫉妒和纷争相反的是谁呢？岂非相反基督吗？因为基督在自己身上医治这些肉身的私欲，但在任何人身上都不喜爱它们。因此，圣教会只要还拥有这样的肢体，就还没有瑕疵或皱纹。此外，还有其他的罪过，为此整个教会每日呼求说：\BibleQuote{宽恕我们的罪债：}并且，免得我们以为属神的人可免于这些，不仅任何属血气的人，就连属神的人中也没有任何人，唯有那倚靠在主怀里的，以及祂特别爱的那一位说：\BibleQuote{如果我们说我们没有罪，就是自欺，真理也不在我们内。}但在每一种罪中，在较大的罪中更多，在较小的罪中较少，总有一种相反义德的私欲。关于基督，经上记载：\BibleQuote{祂由天主成了我们的智慧、正义、圣化者和救赎者。}因此，在每一种罪中无疑都有相反基督的私欲。但当那\BibleQuote{医治我们一切疾病}的主引领教会到达所应许的病愈时，在它的任何肢体中都不会再有丝毫的瑕疵或皱纹。那时，肉身绝不会再相反圣神，因此圣神也不必再相反肉身。那时这一切争斗将止息，那时两种实体将有最高的和谐；那时将没有人是属血气的，甚至肉身本身也将是属神的。因此，每个跟随基督生活的人如何对待自己的肉身，既相反其邪恶的私欲（他予以约束，以待将来治愈；他予以抑制，尚未痊愈），又滋养和珍视其良善的本性，因为\BibleQuote{从来没有人憎恨过自己的肉身}，基督对教会也是如此，只要把更大的事与较小的事相比是允许的。因为祂既以责斥压制它，免得它骄傲膨胀而不可收拾；又以安慰扶助它，免得它因软弱沉重而跌倒。因此，宗徒说：\BibleQuote{如果我们先省察自己，就不致受审判；但我们受审判时，是受主的惩戒，免得我们和这世界一同被定罪。}在圣咏中：\BibleQuote{我心内的忧愁虽多，你的安慰却使我喜乐。}因此，我们那时该期望肉身完全康复而无任何反对，正如基督的教会那时必然稳妥而无任何恐惧。\par

26. 以上所述对于为真正的节制辩白，反对那些虚伪节制的摩尼教徒，已经足够；以免当节制约束并抑制我们最低的部分，即身体，使之远离过度与非法的快乐时，这丰饶而荣耀的节制之功被相信不是有益地管教，而是敌对地迫害。确实，身体虽与灵魂的本性不同，却并非与人的本性相异：因为灵魂并非由身体组成，但人却由灵魂和身体组成；而且，天主所释放的，祂释放整个人。因此，我们的救主自己也承担了整个人，甘愿释放祂所创造的我们身上的一切。凡持有与此真理相反观点的人，克制情欲对他们有何益处？即使他们确实克制了一些情欲。通过节制，他们身上有什么能被洁净呢？既然他们的这种节制本身是不洁的？它甚至不应被称为节制。因为持有他们的观点是魔鬼的毒药；而节制是天主的恩赐。但是，正如并非每一个忍受任何事物、或以最大的忍耐忍受任何痛苦的人，都拥有那同样为天主所赐、称为忍耐的德行；因为许多人忍受许多折磨，是为了不背叛他们罪恶同谋者或他们自己；许多人是为了满足炽热的情欲，并获取或不放弃那些被邪恶之爱链条束缚的事物；许多人为不同且有害的错误所坚持；对所有这些，我们绝不能说他们拥有真正的忍耐：同样，并非每一个在任何事上克制、或甚至惊人地抑制肉体或心灵的情欲的人，都可以说拥有我们正在讨论其益处与美善的节制。因为有些人——说来也许惊人——通过无节制来约束自己：好比一个女人因向奸夫发誓而克制不与丈夫同房。有些人通过不义来约束：好比配偶不向对方交出性生活的债务，因为他或她已经能够克服这种身体的欲望。也有些人是因虚假的信仰而受骗，并怀着虚妄的希望，追随虚妄的事物，从而约束自己：其中包括所有异端者，以及任何在宗教名义下被任何错误蒙蔽的人；他们的节制如果是真的，他们的信德也应该是真的；但由于那不能称为信德，因它是虚假的，毫无疑问那也不配称为节制。怎么？我们准备称那种节制为罪吗？而我们必须真实地说节制是天主的恩赐。愿如此可憎的疯狂远离我们的心。但蒙福的宗徒说：\BibleQuote{凡不出于信德的都是罪。}因此，凡没有信德的，就不应称为节制。\par

27. 也有一些人，在公开事奉邪魔时，克制身体的快乐，以便通过它们来满足非法的快乐，而他们克制不了这种快乐的暴力和炽热。因此，（举一个例子，因篇幅原因略去其余）有些人不靠近自己的妻子，却仿佛清洁，试图通过邪术去接触别人的妻子。啊，奇妙的节制，不如说是独特的邪恶与不洁！因为，如果是真正的节制，肉体的情欲更应该克制奸淫，而不是为了犯奸淫而克制婚姻。的确，婚姻的节制通常缓解肉体的情欲，并约束其缰绳，但程度是：即使在婚姻本身中，它也不以过度放纵而奔放，而是要遵守一种尺度，或是由于配偶的软弱而应有的尺度（宗徒对此不是命令，而是允许），或是适于生育子女的尺度（这曾是古时圣洁的父母唯一的性交理由）。然而，节制这样做，即调节并以某种方式限制已婚者肉体的情欲，并在固定范围内秩序化其不安和失调的动向，它善用了人的恶，那人因它而被造，且它愿意使之成为完美的善：正如天主为着祂使之成全于善的人的缘故，也善用恶人。\par

28. 因此，我们绝不能说那节制（圣经论及它时说：\BibleQuote{这本身就是智慧，知道它的恩赐是谁的}）也为那些通过节制要么服务于错误，要么克服较小欲望以便实现其他被更大欲望所克服的人所拥有。但真正的节制从天而降，不愿以恶压恶，而是以善治愈一切恶。简言之，要概括其行动方式，节制的作用是警戒并克制和治愈一切与智慧之乐相对立的情欲之乐。无疑，那些将节制仅限于克制身体情欲的人，把它限制得太狭窄了；那些说节制在于统治一般欲念或贪欲的人，说得更好。这种欲念被定为过错，不仅关乎身体，也关乎灵魂。因为，如果说身体的欲念在于奸淫和醉酒，那么强烈的敌意、争执、嫉妒、最后甚至仇恨，难道它们的实行是在身体的快乐中，而不是在灵魂的动荡和烦乱中吗？但宗徒称这些全部为\BibleQuote{本性的工作}，无论它们属于灵魂还是专属于肉体，因为他以肉体的名称称呼人本身。因为，凡不称为天主之工的，就是人之工；因为做这些事的人，是按自己生活，而不是按天主生活，就他做这些事而言。但还有其他人的工作，更应称为天主之工。因为宗徒说：\BibleQuote{是天主在你们内推动，使你们愿意并实行祂的善意。}由此也有一句：\BibleQuote{凡受天主圣神引导的，都是天主的子女。}\par

29. 如此，人的精神隶属于天主圣神，就与肉体相争，即与自身相争：但为自身，为的是那些或出于肉体、或出于灵魂、依人而不依天主的冲动——这些冲动仍因人所患的病而存在——能藉节制受到约束，从而获得健康；人不再依人而生活，如今能说：\BibleQuote{但是，我生活，已不是我生活，而是基督在我内生活。}因为那\Quote{不是我的}之处，正是我更幸福之处；当任何依人的邪恶冲动出现时，凡以心神遵守天主的法律之人不予同意，让他也说这话：\BibleQuote{现在做这事的不再是我。}那些话正是向这样的人说的，我们作为他们的同伴与分享者，应当听从：\BibleQuote{既然你们与基督一同复活了，就该追求天上的事，在那里基督坐在天主的右边；你们该思念天上的事，不该思念地上的事，因为你们已经死了，你们的生命已与基督一同藏在天主内；当基督、你们的生命显现时，你们也要与祂一同出现在光荣中。}让我们明白他在对谁说话——更确切地说，让我们更专注地聆听。还有什么比这更明白？还有什么比这更清楚？他肯定是在对那些已经与基督一同复活的人说话——当然不是指肉体，而是指精神：他称他们为死人，而因此他们更活着；因为他说：\BibleQuote{你们的生命已与基督一同藏在天主内。}这样死去的人的话是：\BibleQuote{但是，我生活，已不是我生活，而是基督在我内生活。}因此，那些生命已藏在天主内的人，被劝勉和鼓励去杀死他们属于地上的肢体。因为紧接着说：\BibleQuote{所以，你们要杀死你们的肢体，那属于地上的。}恐怕有人因过于迟钝以为要杀死那可见的身体的肢体，他立刻解释他所说的是什么：\BibleQuote{淫乱、污秽、邪情、恶欲和贪婪，贪婪即偶像崇拜。}但难道能相信，那些已经死了、生命已与基督一同藏在天主内的人，仍在行淫乱？仍生活在不洁的习惯和行为中？仍受邪情、恶欲和贪婪的奴役？哪个疯子会这样想这些人？那么，他愿意他们杀死什么呢？那就是冲动本身——它们仍凭自己的某种侵入而活着，没有我们心神的同意，没有身体肢体的行动。它们又如何藉节制的作为而被杀死呢？除非我们以心神不同意它们，也不把身体的肢体交给它们作武器；而且，更重要的是，要以更警觉的节制来留意：我们自己的思想本身，虽然以某种方式被它们的暗示所触动，并好像低语，但却从它们那里转开，不从它们那里接受快乐，而转向更令人愉悦的关于天上之事的思念。正因为如此，在讲话中提到它们，是为了让人不留在它们里面，而是逃离它们。而这得以实现，如果我们有效地聆听，依赖祂的帮助——祂通过祂的宗徒颁布这道命令：\BibleQuote{你们该追求天上的事，在那里基督坐在天主的右边。你们该思念天上的事，不该思念地上的事。}\par

30. 但是，在他提到这些邪恶之后，他又加上说：\BibleQuote{因这些事，天主的义怒临到悖逆之子的身上。}这确实是一个有益警告，使信徒不致以为仅凭信德就能得救，即使他们仍生活在这些邪恶中；宗徒雅各伯以最清晰的言论反对那观念，说：\BibleQuote{若有人说他有信德，却没有行为，这信德能救他吗？}因此，外邦人的导师在这里也说了，因这些邪恶，天主的义怒临到悖逆之子的身上。但当他说：\BibleQuote{你们从前也生活在其中，那时你们是活在这些事中的；}这充分表明他们现在不再生活在其中了。的确，他们已经向这些事死了，好使他们的生命与基督一同藏在天主内。既然他们不再活在这些事中，他们现在被命令杀死这些事。的确，他们自己不再活在其中，但这些事仍然活着——正如我上面已经稍微指出的——并且被称为他们的肢体，即那些住在他们肢体中的过犯；这是以一种修辞方式，以包含者指被包含者；正如所说：\Quote{整个论坛都在谈论它}，是指论坛中的人谈论它。正是以这种修辞方式，在圣咏中唱道：\BibleQuote{愿整个大地敬拜祢}——即在大地上的所有人。\par

31. \Quote{但现在你们也要，}他说，\Quote{除去这一切；}并且他还列举了几种类似的恶事。但为什么他说\Quote{你们除去这一切}还不够，还要加上连词说\Quote{你们也要？}除非是恐怕他们以为，因为他们之所以行这些恶事并生活在其中而免受惩罚，是因为他们的信德使他们脱离了那临到不信之子的忿怒——他们行这些事并生活在其中而没有信德。\Quote{你们也要，}他说，\Quote{除去这些恶事，}因为天主的忿怒因这些事临到不信之子；也不要因信德的功劳而自诩免于这些事的惩罚。但他不会对那些已经除去恶事、不再同意这些过失、也不将他们的肢体作为不义的武器献给罪恶的人说\Quote{你们除去，}除非是因为圣人的生命在于这种已实际完成的行为，并且只要我们还属于必朽的，就仍要继续从事这项工作。因为，只要圣神与肉身相争，这事便以极大的热诚进行着，藉着圣德的甘饴、贞洁的爱、属灵的力量和节制之美，抵制邪恶的快乐、不洁的情欲、肉体和可耻的冲动；这样，它们就被那些对它们已经死去、不因同意而生活在其中的人所除去。我说，它们就是这样被除去的，同时它们被持续的节制所压制，以免再次兴起。无论谁，若自以为安全而停止这种除去，它们就会立刻攻击心灵的堡垒，并把它从那里除去，使之成为它们的奴隶，以卑劣和可耻的方式将其俘虏。那时，罪将在人的必朽身体中为王，使人顺从它的欲望；那时，人将把肢体作为不义的武器献给罪；那人的末后境况比先前更糟。因为未曾开始这种战斗远比开始后放弃战斗、从好战士甚至得胜者变成俘虏更能容忍。因此，主不是说\Quote{谁开始，}而是\Quote{谁坚持到底，才可得救。}\par

32. 但无论是激烈争战以免被征服，还是屡次得胜，甚至是意外而轻松的得胜，让我们将荣耀归于那赐予我们节制的天主。让我们记住，一位义人曾说：\Quote{我永不动摇；}并向他显示，他这样说何等轻率，将从上所赐的恩宠归功于自己的力量。但这是我们从他的忏悔学到的：因为不久后他加上说：\Quote{主，祢的旨意赐予了我的美丽以力量；但祢转面不顾，我就惊慌失措。}藉着治疗的眷顾，他暂时被他的统治者遗弃，以免他因致命的骄傲而遗弃他的统治者。因此，无论是此时我们与过失搏斗以制服它们并使之减弱，还是将来——最终之时——我们将没有一切仇敌，因为一切感染都已除去，我们如此被对待是为了我们的健康，好使\Quote{凡夸耀的，应因主而夸耀。}\par

\SourceRecord{Translated by C.L. Cornish.；Nicene and Post-Nicene Fathers, First Series；Vol. 3.；Edited by Philip Schaff.；Buffalo, NY: Christian Literature Publishing Co.,；1887.；<http://www.newadvent.org/fathers/1308.htm>.}
